% Chapter 10 of The Hydrogen Atoms of Deep Learning: grokking and the Fourier circuit.
\newcommand{\grokZp}{\Z_p}
\newcommand{\grokmem}{t_{\mathrm{mem}}}
\newcommand{\grokt}{t_{\mathrm{grok}}}
\newcommand{\grokcw}{\bar\gamma}
\newcommand{\grokii}{\mathrm{i}}

\casechapter{Grokking}{Modular addition and its Fourier circuit}{Modular addition, quadratic two-layer network}{ch:grokking}
\begin{refsection}[chapters/10-grokking.bib]

\begin{chaptersummary}
Small networks trained on part of the addition table of $\Z_p$ fit their training data early and predict the rest of the table only much later: this is grokking. For a two-layer network with quadratic activation the generalizing solution can be written down. Each neuron carries one Fourier frequency, and trigonometric identities turn the sum of the neurons into clocks $\cos(\omega(a+b-c))$ that peak at the correct answer. We prove that margin maximization forces this circuit: every neuron has a single frequency, its readout phase is the sum of its input phases, every frequency is present, and each frequency needs at least six neurons. On a diagonal linear network, weight decay acts exactly as an initialization scale shrinking like $e^{-\lambda t}$, which predicts a delay $\lambda^{-1}\log(\alpha/\alpha_c)$ that simulations confirm to within one percent. The model says nothing about transformers that use a handful of frequencies or a different algorithm, about grokking without weight decay, or about networks trained far from the margin limit.
\end{chaptersummary}

\section{The phenomenon: memorize now, generalize later}\label{sec:grokking:phenomenon}

In 2021 a group at OpenAI trained small transformers on tables of binary operations: addition, subtraction and division modulo $97$, polynomials such as $x^2+xy+y^2$, and composition of permutations in $S_5$. Each network saw a random part of a table and had to fill in the rest \parencite{PowerEtAl2022}.\sidenote{The work was presented at an ICLR 2021 workshop on mathematical reasoning and posted to arXiv in January 2022. The verb \emph{grok}, to understand something completely, comes from Robert Heinlein's novel \emph{Stranger in a Strange Land} (1961).} The networks were two-layer decoder-only transformers of width $128$ with four attention heads and about $4\times10^5$ parameters, trained with AdamW at learning rate $10^{-3}$ and weight decay $1$. On division modulo $97$ with half of the table for training, training accuracy was nearly perfect after about $10^3$ steps, while validation accuracy stayed near chance and approached $100\%$ only after close to $10^6$ steps. Power and coauthors called this delayed generalization \term{grokking}\index{grokking}. They also found that the delay grows steeply as the training set shrinks: at fractions of $25$ to $30\%$ of the table, removing $1\%$ of the data raised the median time to generalization by $40$ to $50\%$, and some operations never generalized. Among the interventions they tried, weight decay improved data efficiency the most.

\begin{observation}[\textcite{PowerEtAl2022}; \textcite{NandaEtAl2023}]\label[observation]{obs:grokking:delay}
A small network trained with AdamW and weight decay on a random fraction of the addition table of $\Z_p$ (between about $25\%$ and $50\%$ of it, for $p$ near $100$) reaches $100\%$ training accuracy early. It stays near chance on the rest of the table for one to three orders of magnitude longer, and then its test accuracy rises to about $100\%$ within a comparatively short interval. The delay grows steeply as the training fraction decreases, and weight decay shortens it.
\end{observation}

What does the network compute once it generalizes? \Textcite{NandaEtAl2023} reverse-engineered a one-layer transformer trained on $30\%$ of the addition table modulo $113$, with full-batch AdamW and weight decay $1$. Its embedding of a number $a$ is concentrated on five ``key frequencies'' $\omega_k=2\pi k/113$, $k\in\{14,35,41,42,52\}$: the network represents $a$ by $\cos(\omega_ka)$ and $\sin(\omega_ka)$. Attention and the MLP combine these into $\cos(\omega_k(a+b))$ and $\sin(\omega_k(a+b))$ by the addition formulas, and the unembedding produces logits close to $\sum_k\cos(\omega_k(a+b-c))$, which are largest at $c=a+b$. Nanda and coauthors also defined progress measures, losses computed after keeping or removing the key-frequency components of the logits, and found three phases: memorization (epochs $0$ to $1.4$ thousand), circuit formation (to $9.4$ thousand), during which the Fourier circuit grows while test loss stays flat, and cleanup (to $14$ thousand), during which weight decay removes the memorizing components and test loss drops. The jump in test accuracy is the visible end of a gradual process, the same kind of hidden progress as in \cref{ch:parity}.

How robust is this? Grokking on algorithmic tasks is reproducible across transformers and multilayer perceptrons and across group operations. \Textcite{LiuEtAl2022} mapped phase diagrams over learning rate and weight decay with four phases (comprehension, grokking, memorization, confusion) and read grokking as a symptom of hyperparameters that sit between fast comprehension and pure memorization. \Textcite{LiuEtAl2023} induced grokking on MNIST, with $1000$ training images and an initialization scaled up nine-fold, and found the delay roughly proportional to the inverse weight decay; on modular addition, a smaller initialization with the weight norm held fixed almost removed it. Delayed generalization has since been reported for convolutional networks on CIFAR-10 \parencite{HumayunEtAl2024}, Gaussian processes and linear regression \parencite{MillerEtAl2024}, and kernel machines with learned features \parencite{MallinarEtAl2025}. Two conclusions are safe. Grokking is a property of a training configuration (initialization scale, weight decay, data fraction) as much as of a task. And on group operations the generalizing networks compute with characters of the group.

These observations raise two questions, and the model of this chapter answers both in a minimal setting. \emph{Which} computation does a generalizing network perform, and why that one rather than another of the many functions that fit the training data? And \emph{why so late}: what sets the delay between fitting the training set and generalizing?

\section{The model: a quadratic network on the addition table}\label{sec:grokking:model}

In this chapter $p\ge3$ is an odd prime, the modulus of the arithmetic; this overrides, locally, the global use of $p$ for a number of features. The elements of $\grokZp=\{0,1,\dots,p-1\}$ are written $a$, $b$, $c$, and arithmetic on them is modulo $p$.\sidenote{The proofs use only that $p$ is odd. Primality matters elsewhere: for composite $p$, a frequency that shares a factor with $p$ cannot separate all the classes on its own.}

\paragraph{The task.} The data are the $p^2$ correct triples $(a,b,a+b)$. A training set $D$ is a random set of pairs $(a,b)$, a fraction of the table, labeled by $a+b$; the test error is the error on the remaining pairs. The input is one-hot: the pair $(a,b)$ is the vector $x=(e_a,e_b)\in\R^{2p}$, which says nothing about the arithmetic. Whatever structure the network finds, it must find in the labels.

\paragraph{The network.} A two-layer network with $m$ hidden neurons, quadratic activation, and $p$ outputs, the logits. The first-layer weight vector $w_r\in\R^{2p}$ of neuron $r$ splits into two functions $u_r,u'_r\colon\grokZp\to\R$, so that $\ip{w_r,x}=u_r(a)+u'_r(b)$, and the readout of neuron $r$ is a third function $v_r\colon\grokZp\to\R$. The logit of class $c$ is
\begin{equation}\label{eq:grokking:net}
  f_c(a,b;\theta)=\sum_{r=1}^m v_r(c)\bigl(u_r(a)+u'_r(b)\bigr)^2 .
\end{equation}
The parameters are $\theta=(\theta_1,\dots,\theta_m)$ with $\theta_r=(u_r,u'_r,v_r)\in(\R^p)^3$; we write $\norm{u}^2=\sum_au(a)^2$ and $\norm{\theta_r}^2=\norm{u_r}^2+\norm{u'_r}^2+\norm{v_r}^2$. The network predicts the class with the largest logit. Training minimizes the cross-entropy of the softmax of the logits over $D$ plus the weight decay $\frac\lambda2\norm\theta^2$, by gradient flow in the analysis and by AdamW in the experiments.

\paragraph{Two structural facts.} The network is \term{homogeneous of degree three} neuron by neuron\index{homogeneous network}: replacing $\theta_r$ by $t\theta_r$ with $t>0$ multiplies the contribution of neuron $r$ to every logit by $t^3$. Homogeneity makes margins the right currency, as in \cref{ch:margin}. And the task is invariant under translations: if $(a,b,c)$ is a correct triple, then so is $(a+g,b+h,c+g+h)$ for all $g,h\in\grokZp$. A problem invariant under a commutative group is diagonalized by the characters of the group, here the functions $a\mapsto e^{2\pi\grokii ka/p}$, so translation invariance is the reason to expect Fourier analysis.

\paragraph{What was thrown away.} The networks in which grokking was discovered are transformers: tokens are embedded, attention mixes the two positions, an MLP with ReLU activations follows, and an unembedding produces logits; training uses Adam with decoupled weight decay on a fraction of the table. The model keeps the two ingredients that the solution needs and nothing else. One-hot inputs with free first-layer weights are a learned embedding. The quadratic activation is the simplest nonlinearity that multiplies, since $(x+y)^2=x^2+2xy+y^2$, and the trigonometric identities behind the circuit need exactly one multiplication; it also makes every logit a polynomial in the weights, so that the Fourier analysis is exact rather than approximate. The margin analysis is done on the full table, which describes what a generalizing network must look like but not why training on part of the table produces one; \cref{sec:grokking:wrong} returns to this. For the delay we use a second, even simpler model, the diagonal linear network of \cref{ch:diagonal} with weight decay, in which the passage from the lazy to the rich regime can be computed exactly.

\begin{dictionary}
one-hot $(e_a,e_b)$ and first-layer functions $u_r,u'_r$ & token embeddings of $a$ and $b$, read by an MLP neuron\\
quadratic activation; the cross term $2u_r(a)u'_r(b)$ & the multiplication done by attention and by the MLP nonlinearity\\
frequency $\zeta$ and phases $\varphi,\psi$ of a Fourier neuron & a key frequency, and the cosine and sine directions of the embedding \parencite{NandaEtAl2023}\\
readout $v_r(c)\propto\cos(\omega_\zeta c+\varphi+\psi)$ & rows of the unembedding; logits $\propto\sum_\zeta\cos(\omega_\zeta(a+b-c))$\\
normalized margin, norm $\norm{\theta}_{2,3}$ & the long-time limit of cross-entropy training with small weight decay\\
average margin over all wrong classes & gap between the correct logit and the mean wrong logit\\
initialization scale $\alpha$; effective scale $\alpha e^{-\lambda t}$ & parameter norm, which decays under weight decay after memorization\\
critical scale $\alpha_c$ & the norm below which training leaves the lazy regime\\
\end{dictionary}

\paragraph{What one had to notice.} Three things. First, that a network trained on part of a group's table can only generalize by finding the group structure, and that for a cyclic group the structure is carried by characters, which a network can store as cosines in its embedding. Second, that one multiplication suffices: $\cos(\omega(a+b))$ is a sum of products of a function of $a$ and a function of $b$, and a square supplies products. Third, that among all networks that classify the table correctly, training selects one by a norm, so the question ``which circuit?'' becomes an extremal problem, and the extremal problem is diagonalized by the same characters.

\section{The solution: a Fourier circuit for addition}\label{sec:grokking:circuit}

\paragraph{Fourier analysis on $\grokZp$.} For $u\colon\grokZp\to\R$ let $\hat u(k)=\sum_{a\in\grokZp}u(a)e^{-2\pi\grokii ka/p}$ for $k\in\grokZp$. In this normalization the inversion formula, Plancherel's identity, and the symmetry of real functions read
\[
  u(a)=\frac1p\sum_{k\in\grokZp}\hat u(k)e^{2\pi\grokii ka/p},\qquad \sum_au(a)^2=\frac1p\sum_k\abs{\hat u(k)}^2,\qquad \hat u(-k)=\overline{\hat u(k)}
\]
(\cref{thm:toolkit:fourier-zp}). All three follow from the orthogonality relation
\begin{equation}\label{eq:grokking:orth}
  \sum_{a\in\grokZp}e^{2\pi\grokii qa/p}=p\,\ind\{q\equiv0\pmod p\}\qquad\text{for every integer }q .
\end{equation}
Because $p$ is odd, the nonzero frequencies pair up as $\pm\zeta$ with $\zeta\in\{1,\dots,\frac{p-1}2\}$; we write $\omega_\zeta=2\pi\zeta/p$. A real function whose transform vanishes outside $\pm\zeta$ is $a\mapsto\rho\cos(\omega_\zeta a+\varphi)$ with $\rho=\frac2p\abs{\hat u(\zeta)}$ and $\varphi=\arg\hat u(\zeta)$.

\begin{lemma}[Sum of cosines]\label[lemma]{lem:grokking:cosines}
For every integer $q$, $\displaystyle\sum_{\zeta=1}^{(p-1)/2}\cos\frac{2\pi\zeta q}{p}=\frac{p\,\ind\{q\equiv0\}-1}{2}$.
\end{lemma}

\begin{proof}
In \eqref{eq:grokking:orth} the term $a=0$ equals $1$, and the terms $a=\zeta$ and $a=p-\zeta$ are complex conjugates with real part $\cos(2\pi\zeta q/p)$. Hence \eqref{eq:grokking:orth} reads $1+2\sum_{\zeta=1}^{(p-1)/2}\cos(2\pi\zeta q/p)=p\,\ind\{q\equiv0\}$.
\end{proof}

With $q=a+b-c$ the lemma says that the logits
\begin{equation}\label{eq:grokking:target}
  f_c(a,b)=\sum_{\zeta=1}^{(p-1)/2}\cos\bigl(\omega_\zeta(a+b-c)\bigr)=\frac{p\,\ind\{c=a+b\}-1}{2}
\end{equation}
are $\frac{p-1}2$ on the correct class and $-\frac12$ on every other class. Each term is a \term{clock}\index{clock}: $\cos(\omega_\zeta(a+b-c))$ is largest when the angle $\omega_\zeta c$ equals $\omega_\zeta a+\omega_\zeta b$, so it compares the sum of two rotations with a third. One frequency already peaks only at the correct class, because $p$ is prime; summing all of them makes every wrong class tie. The question is how a sum of squares of first-layer activations produces one clock.

\paragraph{One frequency at a time.} Fix a frequency $\zeta$ and write $\omega=\omega_\zeta$. Call the neuron with $u(a)=\cos(\omega a+\varphi)$, $u'(b)=\cos(\omega b+\psi)$ and $v(c)=\cos(\omega c+\varphi+\psi)$ the \term{Fourier neuron}\index{Fourier neuron} of frequency $\zeta$ and phases $(\varphi,\psi)$, and write
\[
  n_{\varphi,\psi}(a,b,c)=\bigl(\cos(\omega a+\varphi)+\cos(\omega b+\psi)\bigr)^2\cos(\omega c+\varphi+\psi)
\]
for its contribution to the logits. The readout phase is the sum of the input phases; \cref{thm:grokking:neuron} will explain why. It is convenient to use the readout phase $\chi=\varphi+\psi$ and the difference $\delta=\varphi-\psi$ as coordinates; the function $n_{\varphi,\psi}$ depends on $(\varphi,\psi)$ only through $\chi$ and $\delta$ modulo $2\pi$.

\begin{lemma}[Cancellation within one frequency]\label[lemma]{lem:grokking:block}
Let $\beta_1,\dots,\beta_K$ be real, let $(\varphi_r,\psi_r)$ be phases, and let $\chi_r=\varphi_r+\psi_r$ and $\delta_r=\varphi_r-\psi_r$. Then
\begin{equation}\label{eq:grokking:block}
  \sum_{r=1}^K\beta_r\,n_{\varphi_r,\psi_r}(a,b,c)=\frac12\Bigl(\sum_{r=1}^K\beta_r\Bigr)\cos\bigl(\omega(a+b-c)\bigr)\qquad\text{for all }a,b,c\in\grokZp
\end{equation}
if and only if, for all integers $j$ and $s$ with $\abs j\le2$, $\abs s\le1$ and $(j,s)\ne(0,0)$,
\begin{equation}\label{eq:grokking:moments}
  M_{j,s}\defeq\sum_{r=1}^K\beta_re^{\grokii(j\chi_r+s\delta_r)}=0 .
\end{equation}
\end{lemma}

\begin{proof}
The plan is to expand one neuron into cosines of linear forms in $(a,b,c)$ and to read off the coefficient of each. Put $A=\omega a+\varphi$, $B=\omega b+\psi$ and $C=\omega c+\chi$. Using $\cos^2x=\frac12(1+\cos2x)$ and $4\cos A\cos B\cos C=\cos(A+B-C)+\cos(A+B+C)+\cos(A-B+C)+\cos(A-B-C)$,
\begin{align*}
  n_{\varphi,\psi}&=\tfrac12\cos(A+B-C)+\cos C+\tfrac14\bigl[\cos(2A\pm C)+\cos(2B\pm C)\bigr]\\
  &\qquad+\tfrac12\bigl[\cos(A+B+C)+\cos(A-B\pm C)\bigr],
\end{align*}
each $\pm$ standing for the sum of the two terms. Because $\chi=\varphi+\psi$, the first term is $\frac12\cos(\omega(a+b-c))$ for every neuron. Every other term is $\cos(\omega\ell+j\chi+s\delta)$ with a linear form $\ell$ in $(a,b,c)$ and integers $(j,s)$, as follows: $C$ has $(\ell;j,s)=(c;1,0)$; $2A+C$ has $(2a+c;2,1)$; $2A-C$ has $(2a-c;0,1)$; $2B+C$ has $(2b+c;2,-1)$; $2B-C$ has $(2b-c;0,-1)$; $A+B+C$ has $(a+b+c;2,0)$; $A-B+C$ has $(a-b+c;1,1)$; and $A-B-C$ has $(a-b-c;-1,1)$. Up to sign, these pairs $(j,s)$ are the seven nonzero pairs with $\abs j\le2$ and $\abs s\le1$. Weighting by $\beta_r$ and summing over $r$, the term with $(\ell;j,s)$ contributes its coefficient times $\operatorname{Re}\bigl(e^{\grokii\omega\ell}M_{j,s}\bigr)$, and $M_{-j,-s}=\overline{M_{j,s}}$ because the weights are real. If \eqref{eq:grokking:moments} holds, then every such contribution vanishes, and \eqref{eq:grokking:block} follows.

Conversely, the eighteen functions $e^{\pm\grokii\omega(a+b-c)}$ and $e^{\pm\grokii\omega\ell}$ are characters of $\grokZp^3$ with distinct frequency vectors: two of these vectors differ in some coordinate by $1$, $2$ or $4$, which multiplied by $\zeta$ is not divisible by $p$. Characters with distinct frequency vectors are orthogonal by \eqref{eq:grokking:orth}. So if \eqref{eq:grokking:block} holds, then the coefficient of each $e^{\grokii\omega\ell}$ on its left side vanishes, and that coefficient is a nonzero multiple of the corresponding $M_{j,s}$.
\end{proof}

In the language of numerical integration, \eqref{eq:grokking:moments} says that the points $(\chi_r,\delta_r)$ with weights $\beta_r$ integrate every trigonometric polynomial of degree at most $2$ in $\chi$ and at most $1$ in $\delta$ exactly, as the uniform measure of total mass $\sum_r\beta_r$ does. Such a set of weighted points is a \term{cubature rule}\index{cubature rule} on the torus, and every cubature rule of this kind gives a circuit.

\begin{theorem}[A Fourier network for modular addition]\label{thm:grokking:fourier}
For each frequency $\zeta\in\{1,\dots,\frac{p-1}2\}$ and each pair of phases $(\varphi,\psi)\in\{0,\frac{2\pi}3,\frac{4\pi}3\}^2$, take the neuron with $u(a)=\cos(\omega_\zeta a+\varphi)$, $u'(b)=\cos(\omega_\zeta b+\psi)$ and $v(c)=\frac29\cos(\omega_\zeta c+\varphi+\psi)$. The resulting network of width $\frac92(p-1)$ has logits
\[
  f_c(a,b)=\frac{p\,\ind\{c=a+b\}-1}2\qquad\text{for all }a,b,c\in\grokZp .
\]
In particular it classifies the whole table correctly, and every margin $f_{a+b}(a,b)-f_c(a,b)$ with $c\ne a+b$ equals $\frac p2$.
\end{theorem}

\begin{proof}
Fix $\zeta$. Its nine neurons are the Fourier neurons $n_{\varphi,\psi}$ with weights $\beta=\frac29$. To compute the moments, write $j\chi+s\delta=(j+s)\varphi+(j-s)\psi$; the sum over the grid factorizes,
\[
  M_{j,s}=\frac29\Bigl(\sum_{\varphi}e^{\grokii(j+s)\varphi}\Bigr)\Bigl(\sum_{\psi}e^{\grokii(j-s)\psi}\Bigr),
\]
and by \eqref{eq:grokking:orth} with modulus $3$ the first factor vanishes unless $3$ divides $j+s$, the second unless $3$ divides $j-s$. If $3$ divides both, then it divides $2j$ and $2s$, hence $j$ and $s$, and $\abs j\le2$, $\abs s\le1$ force $j=s=0$. So \eqref{eq:grokking:moments} holds, and \cref{lem:grokking:block} shows that the nine neurons of frequency $\zeta$ contribute $\frac12\cdot9\cdot\frac29\cos(\omega_\zeta(a+b-c))=\cos(\omega_\zeta(a+b-c))$. Summing over $\zeta$ gives \eqref{eq:grokking:target}.
\end{proof}

The proof is a bookkeeping of phases. The squares $u(a)^2$ and $u'(b)^2$ create terms that ignore $b$ or $a$; the cross term creates the wanted clock together with the unwanted $\cos(\omega(a-b\pm c)+\cdots)$ and $\cos(\omega(a+b+c)+\cdots)$; and three phases per input make every unwanted term a sum of three unit vectors spaced by $2\pi/3$. We checked the identity numerically for all odd $p\le23$, where it holds to rounding error.

\begin{remark}[Other designs]\label[remark]{rem:grokking:designs}
A pair of Fourier neurons with phases $(\varphi,\psi)$ and $(\varphi,\psi+\pi)$ turns squares into a product, $n_{\varphi,\psi}+n_{\varphi,\psi+\pi}=\bigl[(x+y)^2-(x-y)^2\bigr]z=4xyz$, with $x=\cos(\omega a+\varphi)$, $y=\cos(\omega b+\psi)$ and $z=\cos(\omega c+\varphi+\psi)$. \Textcite{MorwaniEtAl2024} use four such pairs, one for each product in $\cos(A+B-C)=\cos A\cos B\cos C-\sin A\sin B\cos C+\sin A\cos B\sin C+\cos A\sin B\sin C$: eight neurons per frequency. Three pairs suffice: for $\theta\in\{0,\frac\pi3,\frac{2\pi}3\}$ take $(\varphi,\psi)=(\theta,\theta)$ and $(\theta,\theta+\pi)$. Expanding $4xyz$ as in the proof of \cref{lem:grokking:block}, the three unwanted terms carry the phases $4\theta$, $2\theta$ and $-2\theta$, whose exponentials sum to zero over $\theta$, so
\begin{align*}
  \frac13\sum_\theta\bigl(n_{\theta,\theta}+n_{\theta,\theta+\pi}\bigr)&=\frac43\sum_\theta\cos(\omega a+\theta)\cos(\omega b+\theta)\cos(\omega c+2\theta)\\
  &=\cos\bigl(\omega(a+b-c)\bigr).
\end{align*}
These six neurons per frequency are the ``order-6'' block of \textcite{Tian2025}, who assembles global minimizers of a squared loss from such blocks; they give a network of width $3(p-1)$. \Textcite{Gromov2023} gives a construction with independent uniformly random phases, readout weight $2/N$ and $N$ neurons per frequency; there each $M_{j,s}$ is a sum of $N$ independent terms of mean zero, of size about $2/\sqrt N$, so the unwanted terms cancel only as the width grows.
\end{remark}

The six-neuron identity is a three-term decomposition of the trilinear form $\cos(A+B-C)=\operatorname{Re}(e^{\grokii A}e^{\grokii B}e^{-\grokii C})$ into products of cosines, and three is the least possible: it is the number of real multiplications needed to multiply two complex numbers, as in Gauss's trick.\sidenote{Multiplying $x_1+\grokii x_2$ by $y_1+\grokii y_2$ needs the four products $x_jy_k$ in the obvious way, but three suffice: $x_1y_1$, $x_2y_2$, and $(x_1+x_2)(y_1+y_2)$. Over the reals, three is the minimum.} Can a frequency do with fewer than six Fourier neurons? Not if the weights are positive, which is the case that margins will select.

\begin{lemma}[Six neurons per frequency]\label[lemma]{lem:grokking:six}
If $\beta_1,\dots,\beta_K>0$ and the moment conditions \eqref{eq:grokking:moments} hold, then the points $(\chi_r,\delta_r)$, taken modulo $2\pi$, include at least six distinct points. In particular $K\ge6$.
\end{lemma}

\begin{proof}
Neurons with the same point $(\chi_r,\delta_r)$ modulo $2\pi$ contribute the same summand to every moment, so merging them and adding their weights changes neither the hypothesis nor the positivity; assume that the $K$ points are distinct. Consider the six-dimensional space of trigonometric polynomials
\[
  P(\chi,\delta)=\sum_{\abs j\le1}\,\sum_{s\in\{0,1\}}c_{js}e^{\grokii(j\chi+s\delta)} .
\]
If $K<6$, then the $K$ linear conditions $P(\chi_r,\delta_r)=0$ on six unknowns have a nonzero solution $(c_{js})$. Expanding $\abs{P}^2$, every product of two of its exponentials has frequency $(j-j',s-s')$ with $\abs{j-j'}\le2$ and $\abs{s-s'}\le1$, so the moment conditions kill every term with $(j,s)\ne(j',s')$:
\[
  0=\sum_{r=1}^K\beta_r\abs{P(\chi_r,\delta_r)}^2=\sum_{j,s,j',s'}c_{js}\overline{c_{j's'}}\,M_{j-j',s-s'}=M_{0,0}\sum_{j,s}\abs{c_{js}}^2 .
\]
But $M_{0,0}=\sum_r\beta_r>0$ and the coefficients are not all zero, a contradiction.
\end{proof}

The argument is the classical lower bound for cubature rules with positive weights: a rule that is exact on all products $P\bar Q$ of functions from a space of dimension $N$ needs at least $N$ points.\sidenote{The bound is attained by the six-point rule of \cref{rem:grokking:designs}. A rule that attains it is called tight, and a tight rule has equal weights (\cref{ex:grokking:phases}).} Without the positivity the bound fails: a block of four neurons with unequal input amplitudes and readout weights of both signs produces the clock up to a term that depends on $c$ alone, and \textcite{Tian2025} cancels such terms across frequencies to build exact circuits of width $2p$. Those neurons are not Fourier neurons, and the next section shows that margins exclude them.

\section{The solution: margins select the circuit}\label{sec:grokking:margin}

\Cref{thm:grokking:fourier} exhibits a circuit; it does not say why training finds it. \Textcite{MorwaniEtAl2024} proposed that margin maximization alone determines it. Recall from \cref{ch:margin} why margins are the natural target. For a homogeneous network trained by cross-entropy on data that it can classify, the minimizers of the regularized loss converge in direction, as the weight decay tends to zero, to a network that maximizes the normalized margin \parencite{WeiEtAl2019}, and gradient flow on the unregularized loss converges in direction to a Karush--Kuhn--Tucker point of the margin problem \parencite{LyuLi2020}.\sidenote{The first statement is about global minimizers, for the norm that is penalized. The second gives a stationary point of the margin problem, not necessarily a global maximizer.} Weight decay penalizes $\norm\theta^2=\sum_r\norm{\theta_r}^2$. Morwani and coauthors normalize instead by
\[
  \norm{\theta}_{2,3}=\Bigl(\sum_{r=1}^m\norm{\theta_r}^3\Bigr)^{1/3},
\]
the norm matched to neurons that are homogeneous of degree three: with it the margin problem decouples into one problem per neuron and becomes solvable. With the Euclidean norm it does not decouple (\cref{ex:grokking:l2}), and we know of no complete solution in that case. The analysis is on the full table, all $p^2$ pairs.

\begin{definition}[Margins]\label[definition]{def:grokking:margin}
The \term{margin}\index{margin!of a network} of $\theta$ at $(a,b)$ against $c\ne a+b$ is $f_{a+b}(a,b;\theta)-f_c(a,b;\theta)$. The margin $\gamma(\theta)$ of the network is the smallest of these $p^2(p-1)$ numbers, and its \term{average margin}\index{margin!average} $\grokcw(\theta)$ is their average. The \term{normalized margin} of $\theta\ne0$ is $\gamma(\theta)/\norm\theta_{2,3}^3$, which is invariant under $\theta\mapsto t\theta$ for $t>0$. A network of width $m$ is a \term{max-margin network}\index{max-margin network} if it maximizes the normalized margin among networks of width $m$.
\end{definition}

The maximum exists, because $\gamma$ is continuous on the compact set $\{\norm{\theta}_{2,3}=1\}$. The average margin is the class-weighted margin of Morwani and coauthors: the margin averaged over a uniform pair $(a,b)$ and a uniform wrong class. Its virtue is linearity: $\grokcw(\theta)=\sum_r\grokcw(\theta_r)$, the sum of the average margins of the one-neuron networks $\theta_r$.

\begin{lemma}[Average margin of one neuron]\label[lemma]{lem:grokking:cw}
For one neuron $(u,u',v)$, with $\bar v=\frac1p\sum_cv(c)$ and $(a,b)$ uniform on $\grokZp^2$,
\begin{align}
  \grokcw(u,u',v)&=\frac{2p}{p-1}\,\E_{a,b}\bigl[u(a)u'(b)\bigl(v(a+b)-\bar v\bigr)\bigr]\nonumber\\
  &=\frac{2}{(p-1)p^2}\sum_{k\ne0}\hat u(k)\,\hat u'(k)\,\overline{\hat v(k)} .\label{eq:grokking:cw}
\end{align}
\end{lemma}

\begin{proof}
For fixed $(a,b)$ the average of $v(c)$ over the $p-1$ wrong classes is $(p\bar v-v(a+b))/(p-1)$, so the average margin at $(a,b)$ is $\frac p{p-1}\bigl(u(a)+u'(b)\bigr)^2\bigl(v(a+b)-\bar v\bigr)$. Expand the square. The term $u(a)^2$ averages to zero, because for fixed $a$ the class $a+b$ is uniform when $b$ is, so that $\E_b[v(a+b)]=\bar v$; the term $u'(b)^2$ vanishes in the same way. The cross term $2u(a)u'(b)$ remains, which proves the first equality. For the second, insert $v(c)=\frac1p\sum_k\hat v(k)e^{2\pi\grokii kc/p}$ and sum over $a$ and $b$ separately:
\[
  \sum_{a,b}u(a)u'(b)v(a+b)=\frac1p\sum_k\hat v(k)\,\hat u(-k)\,\hat u'(-k).
\]
The term $k=0$ equals $\hat u(0)\hat u'(0)\hat v(0)/p=\sum_{a,b}u(a)u'(b)\bar v$, which the subtraction of $\bar v$ removes. Divide by $p^2$, multiply by $\frac{2p}{p-1}$, replace $k$ by $-k$, and use $\hat v(-k)=\overline{\hat v(k)}$.
\end{proof}

The cancellation of the squares is the reason for averaging uniformly over wrong classes: the terms $u(a)^2v(c)$ and $u'(b)^2v(c)$, which no single neuron can avoid, carry no information about $a+b$ and contribute equally to the correct logit and to the average wrong logit. What remains is a cubic form that the Fourier transform diagonalizes.

\begin{theorem}[The best single neuron]\label{thm:grokking:neuron}
Every neuron satisfies $\grokcw(\theta_r)\le\gamma^\star\norm{\theta_r}^3$, with
\[
  \gamma^\star=\sqrt{\frac2{27}}\,\frac1{\sqrt p\,(p-1)} .
\]
A neuron $\theta_r\ne0$ attains equality if and only if it is a positive multiple of a Fourier neuron: $u(a)=\rho\cos(\omega_\zeta a+\varphi)$, $u'(b)=\rho\cos(\omega_\zeta b+\psi)$ and $v(c)=\rho\cos(\omega_\zeta c+\varphi+\psi)$ for some $\rho>0$, frequency $\zeta\in\{1,\dots,\frac{p-1}2\}$, and phases $\varphi,\psi$.
\end{theorem}

\begin{proof}
The plan is to pair conjugate frequencies, to bound each pair by the arithmetic--geometric mean inequality, and then to use that $s\mapsto s^{3/2}$ grows faster than linearly to put all the weight on one pair. Both sides are homogeneous of degree three, so assume $\norm{\theta_r}=1$.

\emph{Pairing.} Because $u$, $u'$ and $v$ are real, the terms $k$ and $-k$ of \eqref{eq:grokking:cw} are complex conjugates, and $k\ne-k$ because $p$ is odd. With $x_\zeta=\abs{\hat u(\zeta)}$, $y_\zeta=\abs{\hat u'(\zeta)}$ and $z_\zeta=\abs{\hat v(\zeta)}$,
\[
  \grokcw(\theta_r)=\frac4{(p-1)p^2}\sum_{\zeta=1}^{(p-1)/2}\operatorname{Re}\Bigl(\hat u(\zeta)\hat u'(\zeta)\overline{\hat v(\zeta)}\Bigr)\le\frac4{(p-1)p^2}\sum_{\zeta=1}^{(p-1)/2}x_\zeta y_\zeta z_\zeta .
\]
\emph{The constraint.} Plancherel's identity, with the terms $\pm\zeta$ grouped and the term $k=0$ dropped, gives $1=\norm{u}^2+\norm{u'}^2+\norm{v}^2\ge\frac2p\sum_\zeta s_\zeta$ with $s_\zeta=x_\zeta^2+y_\zeta^2+z_\zeta^2$; equality requires $\hat u(0)=\hat u'(0)=\hat v(0)=0$.

\emph{One frequency.} By the arithmetic--geometric mean inequality $x_\zeta y_\zeta z_\zeta\le(s_\zeta/3)^{3/2}$, with equality if and only if $x_\zeta=y_\zeta=z_\zeta$. And $\sum_\zeta s_\zeta^{3/2}\le(\max_\zeta s_\zeta)^{1/2}\sum_\zeta s_\zeta\le\bigl(\sum_\zeta s_\zeta\bigr)^{3/2}$, with equality if and only if at most one $s_\zeta$ is nonzero. With $\sum_\zeta s_\zeta\le p/2$ these give
\[
  \grokcw(\theta_r)\le\frac4{(p-1)p^2}\Bigl(\frac{p}{6}\Bigr)^{3/2}=\sqrt{\frac2{27}}\,\frac1{\sqrt p\,(p-1)} .
\]
\emph{Equality.} Equality holds throughout if and only if the transforms vanish at $0$ and at every frequency except one pair $\pm\zeta$, the three moduli at $\zeta$ are equal (to $\sqrt{p/6}$), and $\hat u(\zeta)\hat u'(\zeta)\overline{\hat v(\zeta)}$ is a positive real, that is, $\arg\hat v(\zeta)=\arg\hat u(\zeta)+\arg\hat u'(\zeta)$. By the remark before \cref{lem:grokking:cosines}, these are exactly the functions $\rho\cos(\omega_\zeta a+\varphi)$, $\rho\cos(\omega_\zeta b+\psi)$ and $\rho\cos(\omega_\zeta c+\varphi+\psi)$ with $\rho=\frac2p\sqrt{p/6}=\sqrt{2/(3p)}$. Conversely, such a neuron makes every inequality of the chain an equality.
\end{proof}

Two features of the optimum deserve their names. The best neuron is \emph{single-frequency}: mixing frequencies wastes norm, because the objective is cubic and the constraint quadratic, and a cubic gains more from concentration than a quadratic loses. And its phases \emph{add}: the readout phase is the sum of the two input phases, so that the cross term $\cos(\omega a+\varphi)\cos(\omega b+\psi)$ contains $\frac12\cos(\omega(a+b)+\chi)$, which the readout $\cos(\omega c+\chi)$ turns into the clock $\frac14\cos(\omega(a+b-c))$. A single neuron cannot be the whole network, though: its logits $(u(a)+u'(b))^2v(c)$ peak at the correct class only for some pairs. The next theorem assembles optimal neurons into a network all of whose margins are equal.

\begin{theorem}[Max-margin networks; \textcite{MorwaniEtAl2024}, with sharp width]\label{thm:grokking:maxmargin}
\leavevmode
\begin{enumerate}
\item Every network of width $m$ satisfies $\gamma(\theta)\le\gamma^\star\norm{\theta}_{2,3}^3$.
\item A network $\theta\ne0$ attains equality if and only if its logits are
\[
  f_c(a,b;\theta)=\gamma^\star\norm\theta_{2,3}^3\Bigl(\ind\{c=a+b\}-\frac1p\Bigr)\qquad\text{for all }a,b,c .
\]
Then every nonzero neuron is a positive multiple of a Fourier neuron, and each frequency $\zeta\in\{1,\dots,\frac{p-1}2\}$ is carried by at least six neurons.
\item Equality is attained at width $m$ if and only if $m\ge3(p-1)$.
\end{enumerate}
\end{theorem}

\begin{proof}
The plan is to compare the minimum margin with the average margin, which decomposes over neurons, and then to read the equality cases in Fourier space; part (iii) rests on the cubature lemmas of \cref{sec:grokking:circuit}.

\emph{The bound.} A minimum is at most an average, so by \cref{thm:grokking:neuron},
\begin{equation}\label{eq:grokking:chain}
  \gamma(\theta)\le\grokcw(\theta)=\sum_{r=1}^m\grokcw(\theta_r)\le\gamma^\star\sum_{r=1}^m\norm{\theta_r}^3=\gamma^\star\norm\theta_{2,3}^3 .
\end{equation}
\emph{Equality.} If the logits have the stated form, then every margin equals $\gamma^\star\norm\theta^3_{2,3}$, and equality holds. Conversely, suppose that equality holds, and write $\gamma'=\gamma(\theta)$. The first inequality in \eqref{eq:grokking:chain} is then an equality, so all margins equal $\gamma'$: for every $(a,b)$ the wrong logits share a common value $g(a,b)$, and $f_c(a,b)=g(a,b)+\gamma'\ind\{c=a+b\}$. The second inequality is an equality term by term, so by \cref{thm:grokking:neuron} every nonzero neuron is a positive multiple of a Fourier neuron. A Fourier readout has mean zero over $c$, so $\sum_cf_c(a,b)=0$, which gives $g=-\gamma'/p$.

\emph{Every frequency, six times.} Group the neurons by frequency. The neurons of frequency $\zeta$ are $\rho_r$ times Fourier neurons with phases $(\varphi_r,\psi_r)$, and they contribute $B_\zeta=\sum_r\rho_r^3\,n_{\varphi_r,\psi_r}$, whose Fourier transform in $c$, for fixed $(a,b)$, lives on $\pm\zeta$. By \cref{lem:grokking:cosines} the logits are $\gamma'(\ind\{c=a+b\}-\frac1p)=\frac{2\gamma'}p\sum_\zeta\cos(\omega_\zeta(a+b-c))$, and the transform in $c$ of this sum at $\pm\zeta$ is that of its $\zeta$ term alone. Comparing transforms in $c$, frequency by frequency,
\[
  B_\zeta(a,b,c)=\frac{2\gamma'}p\cos\bigl(\omega_\zeta(a+b-c)\bigr)\qquad\text{for every }\zeta .
\]
The right side is not zero, so some neuron carries frequency $\zeta$; by \cref{lem:grokking:block}, applied with the positive weights $\beta_r=\rho_r^3$, the moment conditions hold; and by \cref{lem:grokking:six} at least six neurons carry $\zeta$.

\emph{Width.} Six neurons for each of the $\frac{p-1}2$ frequencies make $m\ge3(p-1)$ necessary. For $m\ge3(p-1)$, use the six-neuron blocks of \cref{rem:grokking:designs} at every frequency and set the other neurons to zero. Multiplying $u$ and $u'$ by $t$ and $v$ by $t^{-2}$ does not change a neuron's function, so the neuron $(\cos(\omega a+\varphi),\cos(\omega b+\psi),\frac13\cos(\omega c+\chi))$ may be replaced by the Fourier neuron $3^{-1/3}(\cos(\omega a+\varphi),\cos(\omega b+\psi),\cos(\omega c+\chi))$. The logits are \eqref{eq:grokking:target}, so all margins are equal and the first inequality in \eqref{eq:grokking:chain} is an equality; every neuron is a Fourier neuron, so the second is one as well.
\end{proof}

So the max-margin networks form a large family, one positive cubature rule per frequency with arbitrary rotations and scalings, but they all compute one function, the clock sum \eqref{eq:grokking:target} up to scale, and they share three visible features: single-frequency neurons, a readout phase equal to the sum of the input phases, and all $\frac{p-1}2$ frequencies present. \Textcite{MorwaniEtAl2024} prove parts (i) and (ii) for width $m\ge4(p-1)$, using their eight-neuron blocks; the six-neuron blocks and \cref{lem:grokking:six} make the width threshold exact. The same paper treats sparse parities and finite groups, where irreducible representations replace the characters.

\section{The solution: why the circuit comes late}\label{sec:grokking:delay}

The max-margin theorem describes the end of training. Why is the end reached so long after the training set is fit? Two groups proposed the same answer in 2023 \parencite{LyuEtAl2024,KumarEtAl2024}: training starts in the lazy regime of \cref{ch:lazy}, where the network fits the training set with its initial tangent kernel and memorizes, and reaches the rich regime, where margin-maximizing features form, only later.

\begin{theorem}[\textcite{LyuEtAl2024}; stated without proof]\label{thm:grokking:lyu}
Let $f(\theta;x)$ be $L$-homogeneous and twice continuously differentiable in $\theta$, let $f(\bar\theta;x)=0$ for every input $x$, and run gradient flow on $\Lhat(\theta)+\frac\lambda2\norm\theta^2$ from $\theta(0)=\alpha\bar\theta$, with weight decay $\lambda=\lambda(\alpha)$ of order $\alpha^{-q}$ for a constant $q>0$. Fix a constant $c\in(0,1)$ and let $\alpha\to\infty$.
\begin{enumerate}
\item At time $(1-c)\lambda^{-1}\log\alpha$ the predictor converges pointwise, after division by a normalizing factor in the case of classification, to a kernel predictor for the tangent features $\nabla_\theta f(\bar\theta;x)$: the maximum-$\ell_2$-margin classifier, if these features separate the training data, or the minimum-norm interpolant, for regression with linearly independent features.
\item At suitable times between $\lambda^{-1}\log\alpha$ and $(1+c)\lambda^{-1}\log\alpha$, every limit point of $\theta/\norm\theta$ (classification) or of $\theta$ (regression) is a Karush--Kuhn--Tucker point of $\min\frac12\norm\theta^2$ subject to $y_if(\theta;x_i)\ge1$ for all $i$, respectively $f(\theta;x_i)=y_i$ for all $i$.
\end{enumerate}
\end{theorem}

For diagonal linear networks the first problem is the maximum-$\ell_2$-margin classifier and the second the maximum-$\ell_1$-margin classifier, which needs far fewer samples for a sparse target; that gap is grokking. \Textcite{KumarEtAl2024} reach the same picture from the other side. They observe grokking in a two-layer network that learns $y=\frac12\ip{w_\star,x}^2$ from Gaussian inputs by plain gradient descent with no weight decay, control it by an output scale that sets the laziness of the dynamics, and find that it needs a large scale and a training set of intermediate size; they derive no formula for the transition time. For modular addition itself, \textcite{MohamadiEtAl2024} prove that permutation-equivariant kernel methods, a model of the lazy phase, need a constant fraction of the table to generalize: the lazy phase can only memorize.

The diagonal network of \cref{ch:diagonal} with weight decay makes the two phases explicit. Let $X\in\R^{n\times d}$ with $n<d$, let $y=Xw_\star$ with a $k$-sparse target $w_\star$, let $w=u\odot u-v\odot v$ with $u(0)=v(0)=\alpha\mathbf1_d$, and run gradient flow on
\[
  \Lhat(w)+\frac\lambda2\bigl(\norm u^2+\norm v^2\bigr),\qquad \Lhat(w)=\frac1{2n}\norm{Xw-y}^2 .
\]
Recall from \cref{ch:diagonal} the hyperbolic-entropy potential $\Phi_\alpha(w)=\frac{\alpha^2}4\sum_iq(w_i/\alpha^2)$ with $q(z)=2-\sqrt{4+z^2}+z\arcsinh(z/2)$, whose gradient is $\nabla\Phi_\alpha(w)_i=\frac14\arcsinh(w_i/2\alpha^2)$. Without weight decay the flow converges to the interpolator $w^\Phi_\alpha=\argmin\{\Phi_\alpha(w):Xw=y\}$, which is close to the minimum-$\ell_2$-norm interpolator when $\alpha$ is large and to the minimum-$\ell_1$-norm interpolator when $\alpha$ is small.

\begin{proposition}[Weight decay shrinks the initialization]\label[proposition]{prop:grokking:decay}
Let $\alpha_t=\alpha e^{-\lambda t}$ and $g(t)=\nabla\Lhat(w(t))$. For every $t\ge0$ and every coordinate $i$:
\begin{enumerate}
\item $u_i(t)v_i(t)=\alpha_t^2$;
\item $\frac{\mathrm d}{\mathrm dt}\nabla\Phi_{\alpha_t}(w(t))=-g(t)$;
\item $w(t)=\argmin\{\Phi_{\alpha_t}(w'):Xw'=Xw(t)\}$.
\end{enumerate}
\end{proposition}

\begin{proof}
The flow is $\dot u=-2u\odot g-\lambda u$ and $\dot v=2v\odot g-\lambda v$. It exists for all time, because the objective decreases along it and is at least $\frac\lambda2(\norm u^2+\norm v^2)$, so the trajectory stays bounded.

\emph{Part (i).} The product rule gives $\frac{\mathrm d}{\mathrm dt}(u_iv_i)=(-2g_i-\lambda)u_iv_i+(2g_i-\lambda)u_iv_i=-2\lambda u_iv_i$, and $u_i(0)v_i(0)=\alpha^2$.

\emph{Part (ii).} Let $\sigma_i=u_i^2+v_i^2$. By part (i), $\sigma_i^2-w_i^2=4u_i^2v_i^2=4\alpha_t^4$, so $\sigma_i=\sqrt{w_i^2+4\alpha_t^4}$, and the flow gives $\dot w_i=2u_i\dot u_i-2v_i\dot v_i=-4\sigma_ig_i-2\lambda w_i$. Let $z_i=w_i/(2\alpha_t^2)$. Since $\frac{\mathrm d}{\mathrm dt}\alpha_t^2=-2\lambda\alpha_t^2$, the quotient rule gives $\dot z_i=(\dot w_i+2\lambda w_i)/(2\alpha_t^2)=-2\sigma_ig_i/\alpha_t^2$, while $\sqrt{1+z_i^2}=\sigma_i/(2\alpha_t^2)$. Hence
\[
  \frac{\mathrm d}{\mathrm dt}\,\frac14\arcsinh z_i=\frac{\dot z_i}{4\sqrt{1+z_i^2}}=-g_i .
\]
\emph{Part (iii).} At $t=0$, $\nabla\Phi_\alpha(w(0))=0$, so by part (ii) $\nabla\Phi_{\alpha_t}(w(t))=-\int_0^tg=-\frac1nX^\top\int_0^t(Xw-y)$ lies in the row space of $X$. The potential $\Phi_{\alpha_t}$ is strictly convex, because $q''>0$, so a point $w'$ minimizes it on the affine set $\{Xw'=Xw(t)\}$ if and only if $\nabla\Phi_{\alpha_t}(w')$ lies in the row space; $w(t)$ is such a point, and the minimizer is unique.
\end{proof}

In words: on a diagonal network, weight decay acts exactly as an initialization scale that shrinks like $e^{-\lambda t}$, and at every moment the network sits at the predictor that the flow of \cref{ch:diagonal} would have chosen from the scale $\alpha_t$, among the predictors with its own training outputs. Nothing in the proof is asymptotic. The rest of the story needs one approximation.

\begin{heuristic}[Heuristic derivation of the grokking time]
Assume that the training residual stays negligible after an initial transient, so that $Xw(t)\approx y$. Then \cref{prop:grokking:decay} gives $w(t)\approx w^\Phi_{\alpha_t}$, and the test error at time $t$ is a fixed function of $\log\alpha-\lambda t$, read off from the static family $a\mapsto w^\Phi_a$.

\emph{Memorization.} At first $\sigma_i\approx2\alpha^2$, so $\dot w\approx-8\alpha^2\nabla\Lhat(w)$, and the residual $Xw-y$ decays at rate at least $8\alpha^2\lmin(XX^\top/n)$. The training set is fit by the time $\grokmem\approx\log(1/\eps)/\bigl(8\alpha^2\lmin(XX^\top/n)\bigr)$, while $\alpha_t\approx\alpha$ and $w\approx w^\Phi_\alpha$ is close to the minimum-norm interpolator, whose test error for a sparse target is near the null risk.

\emph{Generalization.} The static family passes from $\ell_2$-like to $\ell_1$-like behavior when $a$ crosses a critical scale $\alpha_c$, so the test error drops when $\alpha e^{-\lambda t}=\alpha_c$, that is, at
\begin{equation}\label{eq:grokking:tgrok}
  \grokt\approx\frac1\lambda\log\frac\alpha{\alpha_c} .
\end{equation}
\emph{The critical scale.} For $\abs{w_i}\ll a^2$ the potential is quadratic, $\Phi_a(w)\approx\sum_iw_i^2/(16a^2)$; for $\abs{w_i}\gg a^2$ it is $\Phi_a(w)\approx\frac14\sum_i\abs{w_i}(\log(\abs{w_i}/a^2)-1)$. The minimum-norm interpolator $w_2=X^+y$ is spread over all coordinates, with $\norm{w_2}^2\approx\frac nd\norm{w_\star}^2$, while the target is concentrated on $k$ of them, and the minimizer switches from one to the other when their potentials are comparable. For a target with entries $\pm1$ this gives, to logarithmic accuracy, $\norm{w_2}^2/(16a^2)\approx\frac14\norm{w_\star}_1\log(1/a^2)$, that is,
\[
  \alpha_c^2\log\frac1{\alpha_c^2}\approx\frac{\norm{w_2}^2}{4\norm{w_\star}_1}\approx\frac{n}{4d}.\qedhere
\]
\end{heuristic}

The ratio $\grokt/\grokmem$ grows like $\alpha^2\log(\alpha/\alpha_c)/\lambda$: a large initialization and a small weight decay make grokking dramatic. The prediction \eqref{eq:grokking:tgrok} has the form of \cref{thm:grokking:lyu}, with $\log\alpha$ refined to $\log(\alpha/\alpha_c)$; the refinement matters at the moderate $\alpha$ of experiments. The estimate of $\alpha_c$ is cruder than the rest: on random instances it is within a factor of about $1.5$ of the scale at which the static test error halves, an error that enters $\grokt$ only through a logarithm.

\section{Reading the solution}\label{sec:grokking:reading}

The model makes five quantitative predictions.
\begin{enumerate}
\item The logits of every max-margin network are the clock sum $\gamma(\ind\{c=a+b\}-\frac1p)$: the wrong classes tie. As vectors in $\R^p$, the logit vectors of the $p$ possible sums are the vertices of a regular simplex centered at the origin, the geometry of neural collapse (\cref{ch:collapse}).
\item Every neuron of a max-margin network has a single frequency, equal norms $\norm u=\norm{u'}=\norm v$, and phase defect $\chi-\varphi-\psi=0$, where $\varphi$, $\psi$ and $\chi$ are the phases of $u$, $u'$ and $v$ at their common frequency.
\item All $\frac{p-1}2$ frequencies appear, each in at least six neurons, and the maximal normalized margin is $\gamma^\star=\sqrt{2/27}/(\sqrt p(p-1))$, for instance $6.1\times10^{-4}$ for $p=59$.
\item In the diagonal caricature, the training set is fit at $\grokmem\propto\alpha^{-2}$ and the test error drops at $\grokt\approx\lambda^{-1}\log(\alpha/\alpha_c)$; the test error is a function of $\log\alpha-\lambda t$ alone.
\item The delay $\grokt/\grokmem$ grows like $\alpha^2/\lambda$, up to a logarithm.
\end{enumerate}

\begin{figure}[t]
\begin{fullwidth}
\centering
\begin{tikzpicture}
\begin{groupplot}[group style={group size=2 by 1,horizontal sep=17mm},hplot,width=0.475\linewidth,height=0.37\linewidth]
\nextgroupplot[xmode=log,xmin=1e-5,xmax=3000,ymin=-0.03,ymax=1.03,xlabel={time $t$},ylabel={error},
  xtick={1e-4,1e-2,1,1e2},legend style={at={(0.03,0.04)},anchor=south west}]
\addplot[theory3] table[x=t,y=trainrel] {figures/data/10-grokking-diagonal.dat};
\addplot[simulation,mark repeat=3] table[x=t,y=test] {figures/data/10-grokking-diagonal.dat};
\addplot[theory] table[x=t,y=theory] {figures/data/10-grokking-diagonal.dat};
\addplot[guide] coordinates {(429,-0.03) (429,1.03)};
\legend{training error,test error (flow),test error of $w^\Phi_{\alpha e^{-\lambda t}}$}
\nextgroupplot[xmode=log,xmin=0.8,xmax=125,ymin=1.5,ymax=7,xlabel={initialization scale $\alpha$},ylabel={$\lambda\,t_{1/2}$},
  legend style={at={(0.03,0.97)},anchor=north west}]
\addplot[theory,domain=0.9:110,samples=50] {ln(x/0.1369)};
\addplot[simulation] table[x=alpha,y=sim0003] {figures/data/10-grokking-delay.dat};
\addplot[simulation2] table[x=alpha,y=sim001] {figures/data/10-grokking-delay.dat};
\addplot[only marks,mark=triangle*,mark size=2pt,color=sea,mark options={fill=seamist}] table[x=alpha,y=sim003] {figures/data/10-grokking-delay.dat};
\legend{$\log(\alpha/\alpha_c)$,$\lambda=0.003$,$\lambda=0.01$,$\lambda=0.03$}
\end{groupplot}
\end{tikzpicture}
\caption{Grokking in a diagonal linear network with weight decay ($d=100$, $n=30$, a $3$-sparse target with entries $1$, Gaussian design, seed $0$). The flow of $(u,v)$ is integrated with a stiff solver. Left: $\alpha=10$, $\lambda=0.01$. The training error $\norm{Xw-y}^2/\norm y^2$ falls below $10^{-4}$ by $t=0.005$; the relative test error $\norm{w-w_\star}^2/\norm{w_\star}^2$ stays at $0.72$, the value for the minimum-norm interpolator, until $t\approx300$ and falls below $0.01$ by $t\approx690$. The curve is the quasi-static prediction of \cref{prop:grokking:decay}, the test error of $\argmin\{\Phi_{\alpha_t}(w):Xw=y\}$; the dashed line is $\lambda^{-1}\log(\alpha/\alpha_c)=429$, where $\alpha_c=0.137$ is the scale at which that static test error halves. Right: the time $t_{1/2}$ at which the simulated test error halves, for three weight decays and seven initialization scales, against $\log(\alpha/\alpha_c)$; all $21$ points lie within $0.06$ of the line.}\label{fig:grokking:diagonal}
\end{fullwidth}
\end{figure}

\Cref{fig:grokking:diagonal} tests the delay. The left panel is grokking in miniature: the training error falls at $t\approx10^{-3}$, the test error at $t\approx4\times10^2$, and between them the test error sits at the null-like value $0.72$ for five decades of time. After the fast transient the simulated test error is indistinguishable from the quasi-static curve, which is the test error of the static interpolator $w^\Phi_a$ evaluated at $a=\alpha e^{-\lambda t}$, computed without reference to the dynamics. The right panel shows the collapse that the proposition predicts: $\lambda t_{1/2}$ depends on $\alpha$ only through $\log(\alpha/\alpha_c)$, with the same $\alpha_c$ for every $\lambda$. The small upward offset grows with $\lambda$, from $0.005$ at $\lambda=0.003$ to $0.05$ at $\lambda=0.03$; it is the training residual of order $\lambda$ that the quasi-static assumption neglects. The heuristic estimate $\alpha_c^2\log(1/\alpha_c^2)\approx\norm{w_2}^2/(4\norm{w_\star}_1)$ gives $\alpha_c=0.131$ here and $\grokt=434$.

\begin{figure}[t]
\begin{fullwidth}
\centering
\begin{tikzpicture}
\begin{groupplot}[group style={group size=2 by 1,horizontal sep=17mm},hplot,width=0.475\linewidth,height=0.37\linewidth]
\nextgroupplot[xmode=log,xmin=1,xmax=10000,ymin=-0.03,ymax=1.03,xlabel={AdamW step},ylabel={accuracy or concentration},
  legend style={at={(0.03,0.97)},anchor=north west}]
\addplot[theory3] table[x=step,y=train] {figures/data/10-grokking-modadd-curves.dat};
\addplot[simulation,mark repeat=2] table[x=step,y=test] {figures/data/10-grokking-modadd-curves.dat};
\addplot[theory2] table[x=step,y=conc] {figures/data/10-grokking-modadd-curves.dat};
\legend{training accuracy,test accuracy,Fourier concentration}
\nextgroupplot[xmin=0,xmax=30,ymin=-3.3,ymax=3.3,xlabel={frequency $\zeta$ of the neuron},ylabel={phase defect $\chi-\varphi-\psi$},
  ytick={-3.1416,-1.5708,0,1.5708,3.1416},yticklabels={$-\pi$,$-\frac\pi2$,$0$,$\frac\pi2$,$\pi$},legend style={at={(0.03,0.97)},anchor=north west}]
\addplot[guide] coordinates {(0,0) (30,0)};
\addplot[simulation,restrict expr to domain={\thisrow{relnorm}}{0.3:1.1}] table[x=zeta,y=defect] {figures/data/10-grokking-modadd-neurons.dat};
\addplot[only marks,mark=o,mark size=1.4pt,color=muted,restrict expr to domain={\thisrow{relnorm}}{0:0.3}] table[x=zeta,y=defect] {figures/data/10-grokking-modadd-neurons.dat};
\legend{theory,large neurons,small neurons}
\end{groupplot}
\end{tikzpicture}
\caption{The quadratic network \eqref{eq:grokking:net} trained on modular addition: $p=59$, width $m=256$, $30\%$ of the table for training, full-batch AdamW (learning rate $10^{-3}$, $\beta=(0.9,0.98)$, weight decay $1$), cross-entropy, seed $0$. Left: training accuracy reaches $99\%$ at step $177$, test accuracy at step $5960$. The Fourier concentration is the share of each neuron's first-layer energy at its dominant frequency $\pm\zeta$, averaged with weights $\norm{\theta_r}^3$; it rises steadily from $0.10$ before the test accuracy moves. Right: after $10^4$ steps, the phase defect of every neuron at its dominant frequency. The $140$ large neurons (norm above $0.3$ times the largest) have median concentration $0.91$, defects below $0.13$ in absolute value (median $0.015$), and cover all $29$ frequencies, with $3$ to $8$ of them per frequency.}\label{fig:grokking:modadd}
\end{fullwidth}
\end{figure}

\Cref{fig:grokking:modadd} tests the features. The trained network groks, with a delay of a factor of about $34$, and at the end of training its neurons have the form that \cref{thm:grokking:neuron} prescribes: the large ones are concentrated on single frequencies, and their readout phases equal the sum of their input phases to within a few hundredths of a radian, although nothing in training refers to phases. All $29$ frequencies are present, as \cref{thm:grokking:maxmargin} predicts for max-margin networks and as \textcite{MorwaniEtAl2024} found in their own experiments with quadratic and ReLU networks. The Fourier concentration rises long before the test accuracy, a progress measure in the sense of \textcite{NandaEtAl2023}.

Two predictions fail, and they fail informatively. The network is far from the margin limit: at step $10^4$ its smallest margin is $0.22$ times its average margin, whereas a max-margin network has all margins equal, and its normalized margin is $0.145\,\gamma^\star$. Six frequencies are carried by fewer than six large neurons, so the unwanted terms of \cref{lem:grokking:block} do not cancel exactly; the network classifies correctly anyway, because correct classification needs only the largest logit to be right, not the wrong logits to tie. The features of the max-margin solution are in place long before its margin is.

What about the delay in this network? With weight decay $0.5$, $1$ and $2$ the test accuracy reaches $99\%$ after about $15000$, $6000$ and $4000$ steps, so the product of weight decay and delay stays between $6000$ and $8000$ steps, as $\grokt\propto\lambda^{-1}$ suggests. But with weight decay $4$ the network has not generalized after $40000$ steps, and with $0.3$ it needs more than $30000$ (\cref{ex:grokking:wdsweep}). The diagonal model cannot see the upper end, where weight decay competes with fitting, the confusion phase of \textcite{LiuEtAl2022}.

\section{What the model gets wrong}\label{sec:grokking:wrong}

\paragraph{What was thrown away.} Four simplifications, each with a price. \emph{The task} is an abelian group, so the characters are known in advance and diagonalize everything; natural data have no such basis, and the analysis has nothing to say about them. \emph{The architecture} is a two-layer network with a polynomial activation, which makes the Fourier analysis exact; there is no attention, no factorization into embedding and unembedding, no normalization, and no ReLU. \emph{The margin theory} describes networks at the end of training on the full table: it uses the limit of vanishing weight decay, the $\ell_{2,3}$ norm instead of the $\ell_2$ norm that weight decay penalizes, and all $p^2$ pairs as training data. It therefore explains which features a generalizing network uses, but not why training on a fraction of the table generalizes, which is the phenomenon. \emph{The delay model} is a reparameterized linear model on a sparse regression task, whose only ``feature'' is a choice of coordinates; it shows how weight decay moves a network from the lazy to the rich regime, but in a model with no circuit to form.

\paragraph{Where it breaks.} \emph{Different networks find different algorithms.} \Textcite{ZhongEtAl2023} found that transformers with weak attention learn a ``pizza'' algorithm, which averages the embeddings of $a$ and $b$ instead of multiplying them, so that all pairs with the same sum lie on a line through the origin; its logits carry an extra factor $\abs{\cos(\omega(a-b)/2)}$ and are not clocks. Which algorithm appears depends on the architecture, and margin maximization within one architecture does not predict it.\sidenote{Zhong and coauthors interpolate between constant attention, which learns the pizza, and standard attention, which learns the clock, by an attention-rate parameter; the transition between the two is continuous.} \emph{Transformers use few frequencies.} The transformer of \textcite{NandaEtAl2023} uses five key frequencies, and other runs between two and nine, while every max-margin network of \cref{thm:grokking:maxmargin} uses all $56$; our MLP agrees with the theory, the transformers do not. \emph{Grokking happens without weight decay.} \Textcite{KumarEtAl2024} observe it with plain gradient descent, controlled by laziness; \textcite{ThilakEtAl2022} (unrefereed) observe it with Adam and no regularization, at the onset of cyclic instabilities in which the norm of the last layer grows abruptly; and \textcite{PrietoEtAl2025} show that without regularization the logits grow until floating-point absorption makes the softmax saturate and the gradient vanish exactly, halting learning, while the gradient aligns with a direction that scales the logits without changing predictions. In these settings the clock that sets the delay is not $\lambda^{-1}\log(\alpha/\alpha_c)$. \emph{Grokking happens without circuits.} \Textcite{LeviEtAl2024} show that a linear estimator on Gaussian data groks in thresholded accuracy, with a delay set by the Marchenko--Pastur spectrum (\cref{ch:doubledescent}) while the losses evolve smoothly; \textcite{VarmaEtAl2023} (unrefereed) explain grokking by the relative efficiency of memorizing and generalizing circuits, which predicts a critical dataset size and ``ungrokking'' below it, an axis the initialization scale does not capture. \emph{Non-abelian groups.} For finite groups the characters give way to irreducible representations, and \textcite{ChughtaiEtAl2023} found networks composing representation matrices, with the set of representations learned varying from seed to seed; but \textcite{StanderEtAl2024}, reverse-engineering networks trained on $S_5$ and $S_6$, found neurons that detect membership in cosets of subgroups and little evidence that the embeddings store representation matrices. Even the group-theoretic generalization of the clock can misdescribe what a network computes.

\paragraph{Fine structure.} The repairs come in four directions. \emph{Representation theory} replaces Fourier analysis for finite groups: \textcite{MorwaniEtAl2024} characterize max-margin networks for general group operations, and \textcite{ChughtaiEtAl2023} describe the algorithm by matrix products. \emph{Algebraic structure of optima}: \textcite{Tian2025} decomposes the squared loss of quadratic networks on abelian groups frequency by frequency and builds global minimizers by composing small blocks, including the six-neuron block above and narrower mixtures of width $2p$ with non-Fourier neurons, and reports that most gradient-descent solutions match these constructions. \emph{Dynamics and phase transitions}: \textcite{RubinEtAl2024} treat grokking in two-layer networks, including the quadratic network on modular addition, as a first-order phase transition in an adaptive-kernel theory of the Bayesian posterior; \textcite{ZunkovicIlievski2024} compute exact critical exponents in teacher--student models with local rules; and \textcite{BeckEtAl2025} place grokking in logistic regression at the edge of linear separability. \emph{Fourier features elsewhere}: recursive feature machines, kernel machines with learned feature metrics, grok modular arithmetic with block-circulant features that implement the same Fourier algorithm \parencite{MallinarEtAl2025}; GPT-2-XL fine-tuned on addition represents numbers with sparse Fourier features, low frequencies in the MLPs for magnitude and high frequencies in attention for the residue \parencite{ZhouEtAl2024}; and language models represent days of the week and months of the year on circles \parencite{EngelsEtAl2025}, a cross-link to \cref{ch:superposition}. What is open is a dynamical theory: a proof that gradient descent with weight decay on a fraction of the table finds the circuit, with the time it takes and the set of frequencies it selects.

\section{Model card}\label{sec:grokking:card}

\modelcard{A network fits part of a group's table early and generalizes much later; once it generalizes, it computes with characters of the group.}{Two-layer network with quadratic activation and one-hot inputs on $a+b\bmod p$; for the delay, a diagonal linear network with weight decay and large initialization.}{An explicit Fourier circuit; max-$\ell_{2,3}$-margin networks consist of Fourier neurons with additive phases, use every frequency with at least six neurons each, and exist exactly at width $m\ge3(p-1)$; weight decay acts as an initialization scale $\alpha e^{-\lambda t}$.}{Logits $\propto\sum_\zeta\cos(\omega_\zeta(a+b-c))$; $\gamma^\star=\sqrt{2/27}/(\sqrt p(p-1))$; phase defect $0$; $\grokt\approx\lambda^{-1}\log(\alpha/\alpha_c)$ and $\grokmem\propto\alpha^{-2}$.}{Transformers use few frequencies or the pizza algorithm; trained networks are far from max margin; grokking without weight decay or without circuits; coset circuits for permutation groups.}{Representation theory of finite groups; algebraic composition of optima; lazy-to-rich dynamics; phase-transition theories; Fourier features in language models.}

The chapter sits on two threads of the book. The first is implicit bias: the max-margin argument is the one of \cref{ch:margin}, transported to a network that is homogeneous of degree three, and the diagonal model of \cref{ch:diagonal} supplies both the potential and the exact effect of weight decay. The second is lazy versus rich training: the memorization phase is the kernel regime of \cref{ch:lazy}, and grokking is the crossing from that regime to the rich one, which weight decay forces at a predictable time. The search phase of \cref{ch:singleindex} and the hidden progress of \cref{ch:parity} are the same story told with small initialization, where the plateau comes before fitting rather than after.

Fourier analysis on groups links this chapter to the Boolean cube of \cref{ch:parity}, where the characters are parities, and to the circular features of \cref{ch:superposition}. The equal-margin logits of a max-margin network form a regular simplex, the configuration that \cref{ch:collapse} derives for the last layer of a classifier.

\begin{exercises}
\item \consolidation (The normalized margin of the circuit.) (a) Show that the network of \cref{thm:grokking:fourier} has normalized margin $0.39\,\gamma^\star$, to two digits. (b) Rescale each neuron, multiplying $u$ and $u'$ by $t=(2/9)^{1/3}$ and $v$ by $t^{-2}$, and show that the normalized margin becomes exactly $\gamma^\star$. Why does rescaling change the normalized margin but not the logits? \label{ex:grokking:normalized}

\item \consolidation\hard (Six phases suffice, and tight rules have equal weights.) (a) Verify the six-neuron identity of \cref{rem:grokking:designs}, and check that its points $(\chi_r,\delta_r)$ are $(\frac{\pi r}3,\pi r)$, $r=0,\dots,5$, up to order. (b) Suppose that six points with positive weights satisfy \eqref{eq:grokking:moments}. Show that every weight equals $\frac16\sum_r\beta_r$. (c) Conclude that the nine phases of \cref{thm:grokking:fourier} are not minimal, and that six is the least number of Fourier neurons per frequency. \hint{for (b), use the space of \cref{lem:grokking:six}: the evaluation map at the six points is a bijection, so for each point there is a $P_r$ that is $1$ there and $0$ at the others. Compute $\sum_s\beta_s\abs{P_r(\chi_s,\delta_s)}^2$ in two ways, and use Cauchy--Schwarz in the form $1=P_r(\chi_r,\delta_r)\le\sqrt6\,(\sum\abs{c_{js}}^2)^{1/2}$.}\label{ex:grokking:phases}

\item \consolidation (Why the full table.) (a) Repeat \cref{lem:grokking:cw} with the wrong class drawn from a weighting $\tau(a,b)\in\Delta(\grokZp\setminus\{a+b\})$ that is not uniform, and show that the squares $u(a)^2$ and $u'(b)^2$ no longer cancel in general. (b) Repeat it with $(a,b)$ uniform on a training set $D$ instead of on $\grokZp^2$, and show that the squares cancel if $D$ is a union of full rows $\{a\}\times\grokZp$, but not for a typical random $D$. What does this say about using the margin theory to explain generalization from a fraction of the table? \label{ex:grokking:subset}

\item \consolidation (Weight decay and the Lasso.) (a) Show that the minimum of $\norm u^2+\norm v^2$ subject to $u\odot u-v\odot v=w$ is $\norm w_1$, so that the regularized objective of \cref{sec:grokking:delay} has the same minimum value as the Lasso $\Lhat(w)+\frac\lambda2\norm w_1$, and its global minimizers map to Lasso solutions. (b) For gradient descent with step $\eta$, show that $u_iv_i$ is multiplied by $(1-\eta\lambda)^2-4\eta^2g_i^2$ at each step, so part (i) of \cref{prop:grokking:decay} holds up to a relative error of order $\eta^2$ per step. (c) Derive the memorization time of the heuristic derivation and evaluate it for the run of \cref{fig:grokking:diagonal}. \label{ex:grokking:lasso}

\item \consolidation (A simplex of logits.) Let $\theta$ be a max-margin network. (a) Show that its logit vector $f(a,b)\in\R^p$ depends on $(a,b)$ only through $a+b$, and that the $p$ distinct vectors have equal norms, sum to zero, and have pairwise inner products $-\frac1{p-1}$ times their squared norm. (b) Compare with the simplex equiangular tight frame of \cref{ch:collapse}. \label{ex:grokking:simplex}

\item \variation\hard (Multiplication and even moduli.) (a) Let $g$ generate the cyclic group $\grokZp^\times$ of order $p-1$. Show that $c=ab$ in $\grokZp^\times$ if and only if $\log_gc=\log_ga+\log_gb$ in $\Z_{p-1}$, and build a Fourier network for multiplication. (b) The modulus $N=p-1$ is even, so the frequency $N/2$, whose character $(-1)^x$ is the Legendre symbol after the change of variables, is real and has no conjugate partner. Redo \cref{thm:grokking:neuron} on $\Z_N$ and show that the best single neuron sits at frequency $N/2$, with average margin $\frac{2}{3\sqrt3(N-1)\sqrt N}$, a factor $\sqrt2$ above the best paired frequency. (c) Show that a network whose neurons all sit at frequency $N/2$ gives every wrong class of the same parity as $a+b$ a margin of zero, and conclude that the proof of \cref{thm:grokking:maxmargin} fails for even moduli. \label{ex:grokking:even}

\item \variation (A non-abelian group.) Let $G$ be a finite group with irreducible characters $\chi_\rho$ of dimensions $d_\rho$. (a) Use the orthogonality of characters to show that $\sum_\rho d_\rho\chi_\rho(g)=\abs G\,\ind\{g=e\}$. (b) Write the character table of $S_3$ and check (a). (c) Deduce that the logits $f_c(a,b)=\sum_\rho d_\rho\chi_\rho(c^{-1}ab)-1$ give every margin $\abs G$, the analogue of \cref{thm:grokking:fourier}. (d) Write $\chi_\rho(c^{-1}ab)$ as a trace of a product of three matrices, and explain why a network needs products of matrix entries rather than of numbers; compare with the algorithm of \textcite{ChughtaiEtAl2023}. \label{ex:grokking:s3}

\item \variation (The Euclidean norm.) (a) Show that $\min_{t>0}\bigl(t^2(\norm u^2+\norm{u'}^2)+t^{-4}\norm v^2\bigr)=3\cdot2^{-2/3}(\norm u^2+\norm{u'}^2)^{2/3}\norm v^{2/3}$, so that the weight-decay cost of a neuron grows like the power $\frac23$ of the size of its contribution to the logits. (b) Show that under $\norm\theta_2\le1$ the average margin is maximized only by networks with a single nonzero neuron, and that for $p\ge5$ such a network does not have equal margins. Conclude that the proof of \cref{thm:grokking:maxmargin} does not apply to the norm that weight decay penalizes. \label{ex:grokking:l2}

\item \variation (Unbalanced initialization.) In the diagonal network start from $u(0)=\alpha\mathbf1_d$ and $v(0)=\alpha'\mathbf1_d$ with $\alpha\ne\alpha'$. Show that $u_iv_i=\alpha\alpha'e^{-2\lambda t}$, that part (ii) of \cref{prop:grokking:decay} holds with $\alpha_t^2=\alpha\alpha'e^{-2\lambda t}$, and that $w(t)$ minimizes $\Phi_{\alpha_t}(w')-\frac14\log(\alpha/\alpha')\sum_iw'_i$ on $\{Xw'=Xw(t)\}$. Which sign of target coordinates does a larger $\alpha$ favor? \label{ex:grokking:unbalanced}

\item \simulation (Weight decay and the delay in the network.) Run the script of \cref{fig:grokking:modadd} with weight decay $0.3$, $0.5$, $1$, $2$ and $4$ and $40000$ steps, and record the first step at which the test accuracy reaches $99\%$. The reader should observe about $15000$, $6000$ and $4000$ steps for $0.5$, $1$ and $2$, consistent with $\grokt\propto\lambda^{-1}$ up to a slowly varying logarithm; more than $30000$ for $0.3$; and no generalization within $40000$ steps for $4$, where the network memorizes later and weight decay competes with fitting. Plot the weight norm against time in each run and compare with the effective scale $\alpha e^{-\lambda t}$ of the diagonal model. \label{ex:grokking:wdsweep}

\item \simulation (The critical scale.) Using the script of \cref{fig:grokking:diagonal}, (a) set $\lambda=0$ and check that the test error stays at that of $w^\Phi_\alpha$ forever; (b) for $n/d\in\{0.2,0.3,0.4\}$ and several random designs, compute the scale $\alpha_{1/2}$ at which the static test error halves and compare it with the root of $\alpha_c^2\log(1/\alpha_c^2)=\norm{w_2}^2/(4\norm{w_\star}_1)$. The reader should find agreement within a factor of about $1.5$, and an error in $\lambda\grokt$ of about $\log1.5\approx0.4$, small against $\log(\alpha/\alpha_c)$ when $\alpha$ is large. \label{ex:grokking:alphac}

\item \research (Dynamics on a fraction of the table.) Prove that gradient descent with small weight decay, started from a small random initialization and trained on a random fraction of the table, converges to a network whose neurons are close to Fourier neurons, and estimate the time. \Textcite{MohamadiEtAl2024} prove generalization bounds for regularized global minimizers with large $\ell_\infty$ margin, and \textcite{Tian2025} describes the global minimizers of a squared loss; a proof for the dynamics from random initialization, and a prediction of which frequencies are selected and when, are open. \label{ex:grokking:dynamics}

\item \research (The narrowest circuit.) What is the smallest width of a quadratic network \eqref{eq:grokking:net} whose logits equal the clock sum \eqref{eq:grokking:target} exactly? Fourier neurons need $3(p-1)$, and \textcite{Tian2025} reaches width $2p$ with other neurons. Numerical searches suggest that one frequency cannot be represented exactly by four neurons of that frequency, and that with five the error can be made small only with very large weights. A matching lower bound, and the analogous question for the $\ell_2$ max-margin problem of \cref{ex:grokking:l2}, are open. \label{ex:grokking:narrow}
\end{exercises}

\notesandrefs
Grokking was reported by \textcite{PowerEtAl2022}, and its explanation in terms of Fourier features comes from the mechanistic study of \textcite{NandaEtAl2023}. \Textcite{Gromov2023} trained the two-layer quadratic network of this chapter, with a squared loss, observed the same features, and wrote down an approximate solution with random phases. The effective theory of \textcite{LiuEtAl2022} and the weight-norm picture of \textcite{LiuEtAl2023} first tied the delay to hyperparameters. \Textcite{DaviesEtAl2023} (a workshop paper) relate grokking to epoch-wise double descent, and \textcite{DoshiEtAl2024} study it with corrupted labels, where memorizing and generalizing neurons can be told apart.

The margin analysis is that of \textcite{MorwaniEtAl2024}, with the class-weighted margin as the device that linearizes the problem; their paper treats sparse parities and general finite groups as well. The general margin results behind it are those of \textcite{WeiEtAl2019} and \textcite{LyuLi2020}, discussed in \cref{ch:margin}. The eight-neuron blocks are Morwani and coauthors'; the six-neuron block appears in \textcite{Tian2025}, whose algebraic framework also produces the narrower mixtures mentioned in \cref{sec:grokking:circuit}. The nine-phase grid of \cref{thm:grokking:fourier} is the most symmetric design, and the cubature view of \cref{lem:grokking:block,lem:grokking:six} organizes all of them.

For the delay, \textcite{LyuEtAl2024} prove the dichotomy between early kernel and late margin biases for homogeneous networks, with diagonal linear networks and matrix completion as examples, and they also exhibit ``misgrokking'', in which the early bias generalizes and the late one does not. \Textcite{KumarEtAl2024} argue for the lazy-to-rich transition without weight decay, building on the lazy regime of \textcite{ChizatEtAl2019} and the kernel and rich regimes of diagonal networks of \textcite{WoodworthEtAl2020}. \Textcite{MohamadiEtAl2024} prove the kernel lower bound for modular addition and generalization bounds for quadratic networks with large margin. The weight-decay identity of \cref{prop:grokking:decay} is a short computation in the setting of \cref{ch:diagonal}. AdamW is due to \textcite{LoshchilovHutter2019}; its decoupled weight decay multiplies the parameters by $1-\eta\lambda$ at each step, so that $\eta\lambda$ is the decay rate per step in our experiments. A related account explains grokking by the competition of a sparse and a dense subnetwork \parencite{MerrillEtAl2023}.

\paragraph{Translating the literature.} \Textcite{MorwaniEtAl2024} write a neuron as $(u_a+v_b)^2w$: their $u$, $v$, $w$ are our $u$, $u'$, $v$; their frequency $\zeta$ and phases $\theta^*_u$, $\theta^*_v$, $\theta^*_w$ are our $\zeta$, $\varphi$, $\psi$, $\chi$; their class-weighted margin with uniform weighting is our average margin $\grokcw$; their $L_{2,3}$ norm is our $\norm\cdot_{2,3}$; their $\gamma^*$ is our $\gamma^\star$; their Fourier transform is unnormalized, as ours is. \Textcite{NandaEtAl2023} write $w_k=2\pi k/p$ for the key frequencies, our $\omega_\zeta$ with $\zeta=k$. \Textcite{Gromov2023} writes $N$ for the width (our $m$), $n$ and $m$ for the two inputs (our $a$ and $b$), $q$ for the output (our $c$), and $\phi^{(1)}$, $\phi^{(2)}$, $\phi^{(3)}$ for the phases (our $\varphi$, $\psi$, $\chi$). \Textcite{Tian2025} writes $d$ for the order of the group (our $p$), and the order of a block is our number of neurons per frequency. \Textcite{LyuEtAl2024} use our $\alpha$ and $\lambda$, write $\bar\theta_{\mathrm{init}}$ for our $\bar\theta$, and write $T^\pm_c(\alpha)=(1\pm c)\lambda^{-1}\log\alpha$ for the two times of \cref{thm:grokking:lyu}.

\printchapterbib
\end{refsection}
